\subsection{Proofs of the lemmas related to the \texttt{Try-Cover} procedure}\label{sec:app_try_partiton}

To prove Lemma~\ref{lem:lowerbounds}, we first introduce Definition~\ref{def:br_meta_action} for a given a set of meta-actions $\dreg$. This definition associates to each meta-acxxtion $d$ the set of contracts in which the agent's utility, computed employing the empirical distribution over outcomes returned by the \texttt{Action-Oracle} procedure and the cost of a meta-action introduced in Definition~\ref{def:cost_action}, is greater or equal to the utility computed with the same quantities for all the remaining meta-actions in $\dreg$. Formally we have that:
\begin{definition}\label{def:br_meta_action}
Given a set of meta-actions $\mathcal{D}$, we let $\mathcal{P}_{d_i} (\mathcal{D}) \subseteq \mathcal{P}$ be the set defined as follows:
\begin{equation*}
	\mathcal{P}_{d_i} (\mathcal{D}) \coloneqq \left \{ p \in \mathcal{P} \mid \sum_{\omega \in \Omega}\widetilde{F}_{d_i,\omega}p_{\omega} - c({d_i}) \ge  \sum_{\omega \in \Omega}\widetilde{F}_{d_j,\omega}p_{\omega} - c({d_j})  \,\,\, \forall d_j \in \mathcal{D}\right \}.
\end{equation*}
\end{definition}

It is important to notice that we can equivalently formulate Definition~\ref{def:br_meta_action} by means of Definition~\ref{def:hyperplane}. More specifically, given a set of met actions $\dreg$, for each $d_i \in \dreg$ we let $\mathcal{P}_{d_i} (\mathcal{D}) \coloneqq \cap_{j \in \dreg} H_{ij}$ be the intersection of a subset of the halfspaces introduced in Definition~\ref{def:hyperplane}.

As a second step, we introduce two useful lemmas. Lemma~\ref{lem:union_pregions} shows that for any set of meta-actions $\dreg$, the union of the sets $\preg_d(\mathcal{D})$ over all $d \in \mathcal{D}$ is equal to $\mathcal{P}$. On the other hand, Lemma~\ref{lem:inclusion_upperbounds} shows that the set $\preg_d(\mathcal{D})$ is a subset of the upper bounds $\ureg_d$ computed by the \texttt{Try-Cover} procedure. 

\begin{lemma}\label{lem:union_pregions}
	Given a set of meta-actions $\dreg$ it always holds $\cup_{d\in \mathcal{D}} \preg_d(\mathcal{D}) = \mathcal{P}$.
\end{lemma}
\begin{proof}
	
%The lemma follows by observing that since $|\dreg| \ge 1$, for each $p \in \preg$ there always exits a $d \in \dreg$ such that:
%\begin{equation*}
%\sum_{\omega \in %\Omega}\widetilde{F}_{d_i,\omega}p_{\omega} - c({d_i}) \ge  \sum_{\omega \in \Omega}\widetilde{F}_{d_j,\omega}p_{\omega} - c({d_j})  \,\,\, \forall d_j \in \mathcal{D},
%\end{equation*}
%meaning that $p \in \mathcal{P}_{i} (\mathcal{D})$. Thus $\cup_{d_i\in \mathcal{D}} \preg_i(\mathcal{D}) = \mathcal{P}$, concluding the proof.

The lemma follows observing that, for each $p \in \preg$, there always exits a $d_i \in \dreg$ such that:
\begin{equation*}
	\sum_{\omega \in \Omega}\widetilde{F}_{d_i,\omega}p_{\omega} - c({d_i}) \geq \sum_{\omega \in \Omega}\widetilde{F}_{d_j,\omega}p_{\omega} - c({d_j}) \,\,\, \forall d_j \in \mathcal{D}.
\end{equation*}
This is due to the fact that the cardinality of $\dreg$ is always ensured to be greater than or equal to one. Therefore, for each contract $p \in \preg$, there exists a meta-action $d$ such that $p \in \mathcal{P}_d(\dreg)$, thus ensuring that $\cup_{d\in \mathcal{D}} \preg_d(\mathcal{D}) = \mathcal{P}$. This concludes the proof.
\end{proof}


%\begin{lemma}
%	It holds $P_i(\mathcal{D}) = \cap_{j \in\mathcal { D}} H_{i,j}$
%\end{lemma}
%
%\begin{proof}
%	Let $p$ be in $P_i(\mathcal{D})$. Then, there exists an action $a \in d_i$ such that $p \in P_a(D)$.  
%	Consider a $d_j \in D$. By the definition of $P_a(D)$, for all $a_j\in D_j$, it holds  \[\sum_{\omega \in \Omega}F_{a_i,\omega}p_{\omega} - c_{a_i} \ge  \sum_{\omega \in \Omega}F_{a_{j},\omega}p_{\omega} - c_{a_j}.\]
%	This implies $p \in H_{i,j}$  
%\end{proof}



\begin{lemma}\label{lem:inclusion_upperbounds}
	Under the event $\mathcal{E}_\epsilon$, it always holds $\preg_{d}(\mathcal{D}) \subseteq \mathcal{U}_{d}$ for each meta-action $d \in \mathcal{D}$.
\end{lemma} 

\begin{proof}
	To prove the lemma we observe that, for any meta-action $d_i \in \dreg$, the following inclusions hold:
	\begin{equation*}
		\mathcal{U}_{d_i} = \cap_{j \in \mathcal{D}_{d_i}} \widetilde H_{ij}\supset \cap_{j \in \mathcal{D}_{d_i}} H_{ij} \supset \cap_{j \in \mathcal{D}} H_{ij}= \mathcal{P}_{d_i}(\mathcal{D}).
	\end{equation*}
	The first inclusion holds thanks to the definition of the halfspace $\widetilde{H}_{ij}$ and employing Lemma~\ref{lem:diff_cost_hyperplane}, which entails under the event $\mathcal{E}_\epsilon$. The second inclusion holds because $\mathcal{D}_{d_i}$ is a subset of $\mathcal{D}$ for each $d_i \in \dreg$. Finally, the last equality holds because of the definition of $\preg_{d_i}(\mathcal{D})$.
\end{proof}



\lowerbounds*
\begin{proof}
	As a first step we notice that, if Algorithm~\ref{alg:try_partition} returns $\{\mathcal{L}_d\}_{d \in \dreg}$, then the set $\dreg$ has not been updated during its execution and, thus, we must have $\lreg_d=\ureg_d$ for all $d \in \dreg$. In addition, we notice that, under the event $\mathcal{E}_\epsilon$, thanks to Lemma~\ref{lem:union_pregions} and Lemma~\ref{lem:inclusion_upperbounds} the following inclusion holds:
	$$\bigcup_{d_i \in \dreg} \, \ureg_{d_i} \supseteq \bigcup_{d_i \in \dreg} \preg_{d_i}(\dreg) = \preg.$$
	Then, by putting all together we get:
	\begin{equation*}
	\bigcup_{d_i\in \dreg} \lreg_{d_i} = \bigcup_{d_i\in \dreg} \ureg_{d_i} = \bigcup_{d_i\in \dreg} \preg_{i}(\dreg) = \preg,
	\end{equation*}
	concluding the proof.
\end{proof}




\epsilonbregions*
\begin{proof}
	In order to prove the lemma, we rely on the crucial observation that, for any vertex $p \in V(\mathcal{L}_{d_i})$ of the lower bound of a meta-action $d_i \in \mathcal{D}$, the \texttt{Action-Oracle} procedure called by Algorithm~\ref{alg:try_partition} with $p$ as input either returned $d_i$ or another meta-action $d_j \in \mathcal{D}$ such that $p \in \widetilde{H}_{ij}$ with $d_j \in \mathcal{D}_{d_i}$. 
	%
	%In the following, for the ease of notation we denote $a= a^\star(p)$.
	%
	% then the meta-action returned by \texttt{Action-Oracle} is either $d_i$ or $d_j$ with $p^* \in \widetilde{H}_{ij}$. 
	
	First, we consider the case in which the meta-action returned by \texttt{Action-Oracle} in the vertex $p$ is equal to $d_i$.
	%
	In such a case, $a^\star(p) \in A(d_i)$ thanks to Lemma~\ref{lem:boundedactions} and the following holds:
	\begin{align*}
		\sum_{\omega \in \Omega} \widetilde{F}_{d_i,\omega} \, p_{\omega} - {c}({d_i}) & \ge \sum_{\omega \in \Omega} {F}_{a^\star(p) ,\omega} \, p_{\omega} - {c}_{a^\star(p) }  - 9 B \epsilon m n,
	\end{align*}
	by means of Lemma~\ref{lem:utility_closed}. Next, we consider the case in which the meta-action returned by Algorithm~\ref{alg:action_oracle} is equal to $d_j$ with $p $ that belongs to $\widetilde{H}_{ij}$. In such a case the following inequalities hold:
	\begin{align*}
	\sum_{\omega \in \Omega} \widetilde{F}_{d_i,\omega} \, p_{\omega} - {c}(d_i) & \ge \sum_{\omega \in \Omega} \widetilde{F}_{d_j,\omega} \, p_{\omega} - {c}(d_j) - y \nonumber \\
	& \ge \sum_{\omega \in \Omega} {F}_{a^\star(p),\omega} \, p_{\omega} - {c}_{a^\star(p)}  - 9 B \epsilon m n \nonumber - y,
	\end{align*}
	where the first inequality follows by Lemma~\ref{lem:diff_cost_hyperplane} that guarantees that $p$ belongs to $ \widetilde{H}_{ij} \subseteq H^y_{ij}$ with $y=18 B \epsilon m n^2 + 2  n \eta \sqrt{m}$, while the second inequality holds because of Lemma~\ref{lem:utility_closed}, since $a^\star(p) \in A(d_j)$ thanks to Lemma~\ref{lem:boundedactions} .
	
 	Finally, by putting together the inequalities for the two cases considered above and employing Lemma~\ref{lem:utility_closed}, we can conclude that, for every vertex $p \in V(\lreg_{d})$ of the lower bound of a meta-action $d \in \mathcal{D}$, it holds:
 		\begin{align}
 			\sum_{\omega \in \Omega} {F}_{a,\omega} \, p_{\omega} - {c}_{a}
 			& \ge \sum_{\omega \in \Omega} \widetilde{F}_{d,\omega} \, p_{\omega} - {c}(d) - 9 B \epsilon m n   \nonumber   \\
 			& \ge \sum_{\omega \in \Omega} {F}_{a^\star(p),\omega} \, p_{\omega} - {c}_{a^\star(p)}  - \gamma,\label{eq:gamma_br}
 	 \end{align}
 	for each action $a \in A(d)$ by setting $\gamma \coloneqq 27 B \epsilon m n^2 + 2  n \eta \sqrt{m}$.
 	
 	Moreover, by noticing that each lower bound $\lreg_d$ is a convex polytope, we can employ the Carathéodory's theorem to decompose each contract $p \in \mathcal{L}_{d}$ as a convex combination of the vertices of $\mathcal{L}_{d}$. Formally:
 	%
	\begin{equation}\label{eq:pstar}
		\sum_{p' \in V(\mathcal{L}_{d}) } \alpha(p') \, p'_\omega = p_\omega \quad \forall \omega \in \Omega,
	\end{equation}
	%
	where $\alpha(p') \ge 0$ is the weight given to vertex $p' \in V(\mathcal{L}_{d}) $, so that it holds $\sum_{p' \in V(\mathcal{L}_{d})} \alpha(p') = 1$.
	%
	Finally, for every $p \in \mathcal{L}_{d}$ and action $a \in A(d)$ we have:
	%
	\begin{align*}
		\sum_{\omega  \in \Omega }  F_{a^\star(p),\omega}  \, {p}_\omega - c_{a^\star(p)} & =  \sum_{\omega  \in \Omega } F_{a^\star(p),\omega}  \left( \sum_{p' \in V(\mathcal{L}_{d}) } \alpha(p') \, p'_\omega \right) - c_{a^\star(p)} \\
		& = \sum_{p' \in V(\mathcal{L}_{d}) } \alpha(p') \left( \sum_{\omega  \in \Omega }  F_{a^\star(p),\omega} \, p'_\omega - c_{a^\star(p)} \right) \\
		& \le \sum_{p' \in V(\mathcal{L}_{d}) } \alpha(p') \left( \sum_{\omega  \in \Omega }  {F}_{a,\omega} \, p'_\omega - c_{a} + \gamma\right)    \\
		& =  \sum_{\omega  \in \Omega } {F}_{a,\omega}  \left( \sum_{p' \in V(\mathcal{L}_{d}) } \alpha(p') \, p'_\omega \right) - c_{a} + \gamma\\
		& =  \sum_{\omega  \in \Omega }  {F}_{a,\omega} \, p_\omega  - c_{a} + \gamma,
	\end{align*}
	where the first and the last equalities hold thanks of Equation~\eqref{eq:pstar}, while the second and the third equalities hold since $\sum_{p' \in V(\mathcal{L}_{d})} \alpha(p') = 1$. Finally, the inequality holds thanks to Inequality~\eqref{eq:gamma_br} by setting $\gamma \coloneqq 27 B \epsilon m n^2 + 2  n \eta \sqrt{m}$.
\end{proof}



\partitioncomplexity*
\begin{proof}
%As a first step, we prove that Algorithm~\ref{alg:try_partition} terminates in a finite number of rounds.
%%
%By the way in which Algorithm~\ref{alg:try_partition} intersects the upper bounds with the halfspaces computed with the help of the \texttt{Find-HP} procedure (see Line~\ref{line:upper_bounds}), the algorithm terminates with $\mathcal{L}_d = \mathcal{U}_d$ for all $d \in \mathcal{D}$ after a finite number of rounds as long as each halfspace is computed at most once.
%%
%It is easy to see that this is indeed the case.
%%
%Suppose that the halfspace $\widetilde{H}_{ij}$ for the pair of meta-actions $d_i, d_j \in \mathcal{D}$ has already been computed by the algorithm, then the halfspace could be computed again when:
%\begin{itemize}
%	\item[(i)] \texttt{Action-Oracle} returns a meta-action different from $d_i$ for a vertex $p \in V(\ureg_{d_i})$ of the upper bound $\ureg_{d_i}$ and it also holds that $p \notin \widetilde{H}_{ij}$.
%	\item[(ii)] \texttt{Action-Oracle} returns a meta-action different from $d_j$ for a vertex $p \in V(\ureg_{d_j})$ of the upper bound $\ureg_{d_j}$ and $p \notin \widetilde{H}_{j}$.
%\end{itemize}
%
%
%
As a first step, we observe that Algorithm~\ref{alg:try_partition} terminates in a finite number of rounds.
%
%By the way in which Algorithm~\ref{alg:try_partition} intersects the upper bounds with the halfspaces computed with the help of the \texttt{Find-HS} procedure, the algorithm terminates with $\mathcal{L}_d = \mathcal{U}_d$ for all $d \in \mathcal{D}$ after a finite number of rounds as long as each possible halfspace $\widetilde{H}_{ij}$ is computed at most once.
By the way in which Algorithm~\ref{alg:try_partition} intersects the upper bounds with the halfspaces computed with the help of the \texttt{Find-HS} procedure, the algorithm terminates with $\mathcal{L}_d = \mathcal{U}_d$ for all $d \in \mathcal{D}$ after a finite number of rounds as long as, for each meta-action $d_i \in \dreg$, the halfspaces $\widetilde{H}_{ij}$ with $d_j \in \dreg$ are computed at most once.
%
It is easy to see that this is indeed the case.
%
% This is because each one of the approximate separating hyperplanes defining the boundaries of the halfspaces $\widetilde{H}_{ij}$ is computed by the \texttt{Find-HP} procedure at most once.
%
%Specifically, for every pair of meta-actions $d_i, d_j \in \mathcal{D}$, Algorithm~\ref{alg:find_hyperplane} is called to build the halfspace $\widetilde{H}_{ij}$ only when Algorithm~\ref{alg:action_oracle} returns $d_j$ with $d_j \not \in \dreg_i$ for a vertex $p \in V(\ureg_{d_i})$ of the upper bound $\ureg_{d_i}$.
Specifically, for every meta-actions $d_i \in \mathcal{D}$, Algorithm~\ref{alg:find_hyperplane} is called to build the halfspace $\widetilde{H}_{ij}$ only when Algorithm~\ref{alg:action_oracle} returns $d_j$ with $d_j \not \in \dreg_{d_i}$ for a vertex $p \in V(\ureg_{d_i})$ of the upper bound $\ureg_{d_i}$.
%
%If the halfspace has already been computed, then $\mathcal{U}_{d_i} \subseteq \widetilde{H}_{ij}$ by the way in which upper bounds are updated (see Line~\ref{line:upper_bounds}), and, thus, any vertex of  $\mathcal{U}_{d_i}$ also belongs to the halfspace.
%
%As a result, if \texttt{Action-Oracle} called on a vertex of the upper bound $\mathcal{U}_{d_i}$ returns the meta-action $d_j$, then Algorithm~\ref{alg:try_partition} does \emph{not} compute the halfspace again.
%
If the halfspace has already been computed, then $d_j \in \dreg_{d_i}$ by the way in which the set $\dreg_{d_i}$ is updated. As a result, if Algorithm~\ref{alg:action_oracle} called on a vertex of the upper bound $\mathcal{U}_{d_i}$ returns the meta-action $d_j \in \dreg_{d_i} $, then Algorithm~\ref{alg:try_partition} does \emph{not} compute the halfspace again.
%
%A similar reasoning holds for $\widetilde{H}_{ji}$.
%
%Additionally, Algorithm~\ref{alg:try_partition} ensures that the \texttt{Find-HP} procedure always finds well-defined hyperplanes, since, under the event $\mathcal{E}_\epsilon$, Lemma~\ref{lem:nullintersection} guarantees that the binary search performed in Line~\ref{line:hyperplane_loop} of Algorithm~\ref{alg:find_hyperplane} is always conducted between two contracts for which \texttt{Action-Oracle} returned different meta-actions.


For every $d_i \in \mathcal{D}$, the number of vertices of the upper bound $\mathcal{U}_{d_i}$ is at most $\binom{m + n + 1}{m} = \mathcal{O}(m^{n })$, since the halfspaces defining the polytope $\mathcal{U}_{d_i}$ are a subset of the $m+1$ halfspaces defining $ \preg$ and the halfspaces $\widetilde{H}_{ij}$ with $d_j \in \dreg$. The number of halfspaces $\widetilde{H}_{ij}$ is at most $n$ for every meta-action $d_i \in \dreg$.
%This is because, by Lemma~\ref{lem:nullintersection}, it holds $I(d_i) \cap I(d_j) = \varnothing$ for every $d_i \ne d_j \in \mathcal{D}$, and, thus, the halfspace $\widetilde{H}_{ij}$ is never discovered twice during the execution of the algorithm.
Consequently, since each vertex lies at the intersection of at most $m$ linearly independent hyperplanes, the total number of vertices of the upper bound $\mathcal{U}_{d_i}$ is bounded by $\binom{ m + n +1}{m} = \mathcal{O}(m^{n})$, for every $d_i \in \dreg$

We also observe that, for every $d_i \in \dreg$, the while loop in Line~\ref{line:partition_loop3} terminates with at most $V(\ureg_{d_i}) =  \binom{m + n +1}{m}  = \mathcal{O}(m^{n})$ iterations. This is because, during each iteration, either the algorithm finds a new halfspace or it exits from the loop.
%
At each iteration of the loop, the algorithm invokes Algorithm~\ref{alg:action_oracle}, which requires $q$ rounds.
%
Moreover, finding a new halfspace requires a number of rounds of the order of $\mathcal{O}\left(q \log \left(\nicefrac{Bm}{\eta}\right)\right)$, as stated in Lemma~\ref{lem:rounds_hyperplane}. Therefore, the total number of rounds required by the execution of the while loop in Line~\ref{line:partition_loop3} is at most $\mathcal{O} \left( q \left( \log \left(\nicefrac{Bm}{\eta}\right) + \binom{m + n +1}{m} \right) \right)$.

Let us also notice that, during each iteration of the while loop in Line~\ref{line:partition_loop2}, either the algorithm finds a new halfspace or it exits from the loop. This is because, if no halfspace is computed, the algorithm does \emph{not} update the boundaries of $\ureg_{d_i}$, meaning that the meta-action implemented in each vertex of $\lreg_{d_i}$ is either $d_i$ or some $d_j$ belonging to $\dreg_{d_i}$. Moreover, since the number of halfspaces $\widetilde{H}_{ij}$ is bounded by $n$ for each $\ureg_{d_i}$, the while loop in Line~\ref{line:partition_loop2} terminates in at most $n$ steps. As a result, the number of rounds required by the execution of the while loop in Line~\ref{line:partition_loop2} is of the order of $\mathcal{O} \left( n q \left( \log (\nicefrac{Bm}{\eta}) + \binom{m + n +1}{m}  \right) \right)$, being the while loop in Lines~\ref{line:partition_loop3} nested within the one in Line~\ref{line:partition_loop2}.

Finally, we observe that the while loop in Line~\ref{line:partition_loop1} iterates over the set of meta-actions actions $\dreg$, which has cardinality at most $n$. Therefore, the total number of rounds required to execute the entire algorithm is of the order of $\mathcal{O} \left( n^2 q \left( \log \left(\nicefrac{Bm}{\eta}\right) + \binom{ m + n +1}{m} \right) \right)$, which concludes the proof. 
%Finally, since $q = \mathcal{O}\left(\nicefrac{1}{\rho^4} \log(\nicefrac{1}{\delta})\right)$ (see the proof of Theorem~\ref{thm:finalthm}), the proof is concluded.
\end{proof}




